\scanpage{14}
\mspage{-2.30-}

\uline{Proof}: Let $T\handwritten{\subset}G$ be a maximal torus. Let $<x,y>$ be the form on $\handwritten{\unsure{\Gamma|\Pi|\mathscr T}^{\mathbb R}}(T)$ which defines the norm $\Vert\lambda\Vert$. Let $\handwritten{\ell}_1,\ldots,\handwritten{\ell}_n$ be the linear functionals on $\handwritten{\Gamma^{\mathbb R}}(T)$ given by Proposition 2.14. For any $x\in X_k$, let $I(x)$ be the subset $\handwritten{\{}1,\ldots,n\handwritten{\}}$ given by Proposition 2.14. To estimate $\handwritten{\unsure{\nu|v}}^L(x,\handwritten{\delta})$, say $\handwritten{\delta}=\Delta(\lambda)$. There is always some $\handwritten{\gamma}\in G_k$ such that $\handwritten{\gamma}\circ\lambda\circ\handwritten{\gamma}^{-1}$ is a 1-PS of $T$. Then:

\[
\begin{aligned}
\handwritten{\nu}^L(x,\handwritten{\delta})&=\frac{\handwritten{\mu}^{L}(x,\lambda)}{\Vert\lambda\Vert}\\
&=\frac{\handwritten{\mu}^{L}(\sigma(\handwritten{\gamma},x),\handwritten{\gamma}\circ\lambda\circ\handwritten{\gamma}^{-1})}{\Vert\handwritten{\gamma}\circ\lambda\circ\handwritten{\gamma}^{-1}\Vert}\\
&=\max\handwritten{\left\{}\frac{\handwritten{\ell}_i(\handwritten{\overline}{\handwritten{\gamma}\circ\lambda\circ\handwritten{\gamma}^{-1}})}{\sqrt{\langle\handwritten{\overline}{\handwritten{\gamma}\circ\lambda\circ\handwritten{\gamma}^{-1}},\handwritten{\overline}{\handwritten{\gamma}\circ\lambda\circ\handwritten{\gamma}^{-1}}\rangle}}\ \handwritten{\middle|}\ i\in I(\sigma(\handwritten{\gamma},x))\handwritten{\right\}}
\end{aligned}
\]

But, for each $i$, $\handwritten{\left|}\handwritten{\ell}_i(z)\handwritten{\right|}\big/\sqrt{<z,z>}$ is bounded when $z$ varies over all non-zero points of $\handwritten{\Gamma^{\mathbb R}}(T)$. Therefore $\handwritten{\nu}^L(x,\handwritten{\delta})$ is bounded when $x$ and $\handwritten{\delta}$ vary arbitrarily. Now suppose for $x\in S\handwritten{\subset}X_k$, we seek to minimize $\handwritten{\nu}^L(x,\handwritten{\delta})$. According to the above,

\[
\inf_{\substack{\handwritten{\delta}\in\Delta(G)\\x\in S}}\handwritten{\nu}^{L}(x,\handwritten{\delta})
=\inf_{\substack{\handwritten{\gamma}\in G_k\\x\in S}}
\handwritten{\left[}\inf_{z\in\handwritten{\Gamma^{\mathbb Q}}(T)}
\handwritten{\left[}\max_{i\in I(\sigma(\handwritten{\gamma},x))}
\frac{\handwritten{\ell}_i(z)}{\sqrt{<z,z>}}\handwritten{\right]}\handwritten{\right]}
\]

But the first ``inf'' is essentially a minimum over the finite set of subsets $I\handwritten{\subset\{}1,\ldots,n\handwritten{\}}$ which occur as subsets $I(\sigma(\handwritten{\gamma},x))$. Therefore, this ``inf'' will always be attained for some $x$ and $\handwritten{\gamma}$. As for the inner ``inf'', if we replace the variable \inserted{\handwritten{by}} $z\in\handwritten{\Gamma^{\mathbb R}}(T)$, then the ``inf'' will certainly be attained for some $z$ - for the expression to be

% archive sheet number: 13
